We have used names to denote values: constants, parameters of
procedures and counters of ``for'' loops. Until now, once fixed, the value denoted by such a name
never changed. {\em Variables}, on the contrary, make it possible to store intermediate
results of computations and to modify them during the execution of the program.

\section{Variables and assignment}
A variable\index{Variable} is used to keep a value in the memory
of the computer. In \python, a variable is created by a {\em variable
declaration}\index{Variable declaration}, which gives the name of the variable (which is an identifier)
and its initial value, placed after the sign {\idt =}:

\vspace{\neglskip}\mbox{\blc  \blk\idt  b\ =\ \mcl True}\par
\vspace{\neglskip}\mbox{\blc  \blk\idt  i\ =\ 0}\par
\vspace{\neglskip}\mbox{\blc  \blk\idt  r\ =\ 0.0}\par
\vspace{\neglskip}\mbox{\blc  \blk\idt  s\ =\ {\cha \char34\char34}}\par

\noindent This is written exactly like the constant declarations we have used
so far: in \python, a constant is nothing but a variable whose value is
never changed. It is \python\ that decides which part of the memory and which encoding will be used to
represent the value of the variable, according to the type of this value:
{\idt b} is a boolean variable, {\idt i} an integer variable, {\idt r} a rational
variable and {\idt s} a variable containing a character string (here the empty string). In \python, a variable
therefore always has a value, fixed as soon as it is created. This value can be used in an
expression (such as {\idt 2 * i + 1}) simply by writing the name of the variable.
A variable must not be used before it has been created: if the program begins with
{\idt i = i + 1} without first giving an initial value to {\idt i}, execution stops
with the message `{\tt NameError:\ name 'i' is not defined}'.
The {\em assignment\/} \index{Assignment}instruction
makes it possible to change the
value of a variable. For example the instruction:

\vspace{\neglskip}\mbox{\blc  \blk\blk\idt  b\ =\ \mcl False}\par

\noindent assigns the value ``false'' to the boolean variable {\idt b}. In general, an
assignment instruction has the following form:

\index{Assignment}
\begin{center}
\begin{tabular}{@{}|p{12.5cm}|@{}} \hline
{\em Assignment instruction:}\\ \hline \hline
\begin{minipage}[t]{12.5cm}
{\medskip

\vspace{\neglskip}\mbox{\blc  \blk\blk\idt  variableName\ =\ \rm expression%
\footnote{\ The expression can use the old value of the variable.}}\par

\medskip}
\end{minipage}\\ \hline
\end{tabular}
\end{center}

\noindent The symbol\ {\idt =}\ is read ``becomes equal to''. It must not be confused with the
comparison {\idt ==} (``is equal to'') used in the conditions of {\mcl if} and {\mcl while}.
Executing this instruction consists in
evaluating the ``expression'' to compute its value, which becomes that of the variable named
{\idt  variableName}. The type of the variable is that of the value it contains:

\index{Type compatibility}
\begin{center}
\begin{minipage}[t]{11.5cm}
\begin{tabular}{@{}|l|l|l|@{}} \hline
\multicolumn{3}{@{}|l|@{}}{\em Types of values:} \\ \hline \hline
Example of value & Type & Name of the type\\ \hline
{\mcl True}, {\mcl False} & {\idt bool} & boolean \\
{\idt 0}, {\idt 1}, {\idt -5} & {\idt int} & integer \\
{\idt 0.0}, {\idt 1.5} & {\idt float} & rational \\
{\cha \char34\char34}, {\cha \char34{}abc\char34} & {\idt str} & character string \\
\hline
\end{tabular}
\end{minipage}
\end{center}

\noindent  In \python, the type of a variable is not fixed once and for all:
if {\idt i} is an integer variable, the assignment {\idt i = 1.5} is
accepted and {\idt i} becomes a rational variable. This is rarely a good idea, because the
reader of the program gets lost and the error only shows up later, for example in
{\idt fdn(i)}, which requires an integer. It is better to keep the same type for each variable and to make
{\em conversions}\index{Conversion} explicitly: for example
{\idt r = float(i)} or {\idt i = int(r)}.

\section{Saving the state of the robot}

A first use of names is to save the state of the \robot. Suppose for
example that we want to design a procedure to make a drawing that requires various
changes of the grid size. So that calling
this procedure in a program does not modify the grid size used in the
program, this size must be saved
at the beginning of the procedure and restored at the end. It is enough to give it a
name by a {\em local declaration}\index{Local declaration}
{\idt name = expression} placed in the body of the procedure, as in the following example:

{\raggedright
 %\vspace{\baselineskip}
 \vbox to 0pt{\includevspacePostScript[13mm]{Canards}{0.95}{0.95}}%
 \input python/Ducks.tex
 }

\noindent The declarations {\idt lX = valGsX();\ lY = valGsY()}
compute the values of {\idt valGsX()} and {\idt valGsY()} at the beginning of the procedure and
give them the names {\idt lX} and {\idt lY}, which can be used only in the rest of the
body of the procedure {\idt duck}. They are said to be variables {\em local} to the
procedure: they do not exist outside its body.

\section{Counting}
Suppose we want to
count the number of iterations of a ``while'' loop. We use an integer variable:

\mbox{\blc  \blk\idt  i\ =\ 0}\par

\noindent Before executing the ``while'' loop, the variable must be initialized to 0
(it may have been modified since its declaration):

\mbox{\blc  \blk\blk\idt  i\ =\ 0}\par

\noindent After each pass through the loop, the value of the variable must be increased by 1:

\mbox{\blc  \blk\blk\idt  i\ =\ i\ +\ 1}\par

\noindent Suppose that the value of the variable {\idt i} is 5 {\em before} the assignment
\mbox{\idt  i\ =\ i\ +\ 1}. The value of the expression {\idt  i\ +\ 1}
is $5 + 1 = 6$, so that the value of the variable {\idt i} {\em after} the assignment will be
6. Since this instruction is very frequent, \python\ allows it to be abbreviated to
\mbox{\idt  i\ +=\ 1} (``increase {\idt i} by 1''). Likewise \mbox{\idt  i\ -=\ 1} decreases {\idt i}
by 1 and \mbox{\idt  i\ *=\ 2} doubles it.

For example, this can help us determine the size of the window in which the \robot\
moves on the computer screen:
\input python/ComputeRobotWindowSize.tex

\noindent The variables {\idt width} and {\idt height} are declared in the program,
outside the procedures; they are said to be {\em global}\index{Global variable}. A procedure
can use their value without saying anything special. On the other hand, to {\em modify} them, it must
declare so at the beginning of the body of the procedure with {\mcl global}\index{global}\ {\idt width}.
Without this declaration, \python\ would consider that the assignment {\idt width = 0} creates a
new local variable {\idt width}, unrelated to the global variable, which would keep the
value 0 (and {\idt width += 1} without {\idt width = 0} before it would cause the error
`{\tt UnboundLocalError:\ local variable 'width' referenced before assignment}', the local
variable {\idt width} being used before it has received a value).

The operator {\idt +} applied to two character strings puts them end to end, and the
function {\idt str} gives the character string that represents a number (for
example {\idt str(25)} is {\cha \char34{}25\char34}). A string and a number cannot
be put end to end: {\cha \char34{}(\char34}\ {\idt + width} causes the error
`{\tt TypeError:\ can only concatenate str (not "int") to str}'.

\section{Entering a value from the keyboard}
A value can be obtained by typing it on the computer's
keyboard. To
do so, the function {\idt readInt()} (\underline{read} an \underline{int}eger) waits for an integer to be typed and returns
its value. The question asked by the program (written with {\idt write}, see below)
appears in the bottom bar of the \robot's window, and the answer typed on the keyboard
is displayed right after it (it can also be typed in the terminal from which the program was started).
For example we write:

\mbox{\blc  \blk\blk\idt  write({\cha \char34{}n\ =\ \char34})}\par
\mbox{\blc  \blk\blk\idt  n\ =\ readInt()}\par

\noindent to give the variable {\idt n} the value typed. Typing
mistakes can be corrected with the  %touche d'effacement
\framebox[11.mm][l]{$\leftarrow$~~~}\hspace{0.3ex}
backspace key. The entry of the value must be ended by pressing the %retour a la ligne
\framebox[11.mm][l]{$\hookleftarrow$~~~}\hspace{0.3ex}
return key. If what has been typed is not an integer, the program asks for it
again. On a French (AZERTY) keyboard, the digits can be typed with or without the
shift key.

Likewise {\idt readReal()} (\underline{read} a \underline{real}) returns a rational number (of type {\idt float}), {\idt readString()}
(\underline{read} a \underline{string}) returns the character string typed (without apostrophes or quotation marks) and {\idt readChar()} (\underline{read} a \underline{char}acter) returns
the character of the first key pressed in the \robot's window (without waiting for the
return key). In \pascal, a single procedure {\idt readln} made it possible to read a value
of any type. In \python, each reading function returns a value of a well-defined type
(\python's {\idt input} function always returns a character string, which would have to be
converted into an integer with {\idt int}; {\idt readInt} does it for us). On the other hand some
functions, such as {\idt write}, {\idt str} or {\idt abs}, accept parameters of different
types: they are said to be {\em polymorphic}\index{Polymorphic} (whereas other functions, such as
{\idt readInt} or {\idt sqrt}, are {\em monomorphic}\index{Monomorphic} in the sense that their
parameters and their result always have the same type).

The instruction {\idt readln()} waits until the %retour a la ligne
\framebox[11.mm][l]{$\hookleftarrow$~~~}\hspace{0.3ex}
return key is pressed. It is therefore useful to leave time to read what
the program has written before its execution ends.

\section{Displaying the value of an expression}
To display texts in the terminal (and in the bottom bar of the \robot's
window), one
can use the instruction {\idt write(}$t_1${\idt ,} $t_2${\idt ,} \ldots{\idt )}, which writes
$t_1$, $t_2$, \ldots\ one
after the other, on the same line. The instruction
{\idt writeln(}$t_1${\idt ,} $t_2${\idt ,} \ldots{\idt )} does the same and then goes to the next line.
The $t_i$ can be character strings or integer or rational expressions,
whose value is written: if {\idt x} is 3, {\idt writeln(}{\cha \char34{}x = \char34}{\idt , x)}
writes {\tt x = 3}. One can also convert the value of an expression $E$ into a character string oneself
with {\idt str(}$E${\idt )}, as the program below does. For example
this program reads
an integer $n$
and then writes its factorial $n! = (1 \ast 2 \ast \ldots \ast (n - 1) \ast n)$. To ask again
for the value of $n$ as long as it is incorrect, it uses a recursive function {\idt readN}
(see chapter \ref{ch:rec}) which returns the value read:
\input python/Factorial.tex

\noindent The variable {\idt n} of the function {\idt readN} is local to this function: it is
a different variable from the variable {\idt n} of the program, which receives the value returned by
{\idt readN()}. The maximal value of $n$ is set to 20, as in the \ocaml\ version of
this program, where $21!$ exceeds the largest integer. In \python, integers have no limit:
one could compute exactly $100!$, which is written with 158 digits.
\medskip

\needspace{6\baselineskip}\noindent \rule{137mm}{0.1mm}\par\nopagebreak %---------------------------------------------

\noindent\begin{exercices}{\ \\

\tiret Write a program that draws a frame as close as possible to the edge of the window in which
the \robot\ moves.

\tiret Write a program that sketches the drawings to be completed of figure \ref{fig:completer-les-dessins}
(page \pageref{fig:completer-les-dessins}).

\tiret Babylonian numeration \cite{guitel}, which dates from about 2000 BC, is in
base 60. It is therefore a sexagesimal numeration, just like the one we still use
for telling the time. The units, from 1 to 59, are written additively by accumulating ``nails'' (for
1) and ``chevrons'' (for 10), grouped and stacked as shown in the figure
below.\par\nopagebreak
\noindent\begin{minipage}[b][4.40cm]{\textwidth}\vfill\includePostScript{Babyloniens}{0.95}{0.95}\end{minipage}\par

\noindent Beyond 60, the notation is positional, the value of a digit in
the $p^{\rm th}$ position (counted from right to left starting from zero) being
multiplied by $60^{p}$. The absence of sexagesimal units in a middle position (which
we would write as 0) was first marked by a space and then, around the $3^{\rm rd}$
century BC, by ``two slanting nails''. The Babylonians thus took 15 centuries to
invent a notation for
zero, without however understanding its computational value! For example
the number $3657 = 1 \times 60^{2} + 0 \times 60^{1} + 57 \times 60^{0}$
is represented by juxtaposing the digits $1$, $0$ and
$57 = 50 + 7$:

\par\nopagebreak
\noindent\begin{minipage}[b][2.40cm]{\textwidth}\vfill\includePostScript{Babyloniens3657}{0.95}{0.95}\end{minipage}\par

Write a \python\ program that converts
any po\-si\-tive number into Babylonian nu\-mer\-a\-tion.

\tiret Egyptian numeration ($3^{\rm rd}$
millennium BC) is decimal and additive \cite{guitel}. It uses hieroglyphs to
write the
powers of 10 (the vertical stroke for 1, the handle for 10, the spiral for 100, the lotus for 1000,
the finger for 10,000,
the tadpole for 100,000 and the genie for 1,000,000):\par\nopagebreak
\noindent\begin{minipage}[b][3.50cm]{\textwidth}\vfill\includePostScript{ChiffresEgyptiens}{0.95}{0.95}\end{minipage}\par

\noindent The value of a sequence of signs is the sum of the values
of these signs:\par\nopagebreak
\noindent\begin{minipage}[b][3.07cm]{\textwidth}\vfill\includePostScript{1234567Egyptien}{0.95}{0.95}\end{minipage}\par

\noindent Writing goes from left to right or from right to left, the
reading being done in the direction in which the hieroglyphs face.
Identical signs are grouped on one, two or three lines
as follows:\par\nopagebreak
\noindent\begin{minipage}[b][1.89cm]{\textwidth}\vfill\includePostScript{GroupesEgyptiens}{0.95}{0.95}\end{minipage}\par

\noindent  Between 2000 and 1600 BC, a hybrid representation (additive and partially
positional) appeared. A number less than 1000 placed before, above or below
the sign of the genie of the million, the tadpole, the finger or the lotus represents the product of this number and the
value of this hieroglyph. For example, the number 37,280,549 can be decomposed into
$((372 \ast 100{,}000) + 80{,}549)$ and is written:\par\nopagebreak
\noindent\begin{minipage}[b][3.10cm]{\textwidth}\vfill\includePostScript{37280549Egyptien}{0.95}{0.95}\end{minipage}\par

\noindent Write a \python\ program that converts a number into E\-gyp\-tian
nu\-mer\-a\-tion.

\tiret Greek acrophonic numeration \cite{guitel} was used from
the 5th
century BC in Athens. It is an additive decimal numeration using the auxiliary
base 5 to reduce repetitions. The digits greater
than one are written with
letters (or combinations of letters),
each corresponding to the initial (or the abbreviation) of the name of the
corresponding number:\par\nopagebreak
\noindent\begin{minipage}[b][2.40cm]{\textwidth}\vfill\includePostScript{ChiffresGrecques}{0.95}{0.95}\end{minipage}\par

Write a \python\ program that converts a number less than 100,000 into Greek
nu\-mer\-a\-tion.
To represent large numbers, use the acrophonic numeration
for currencies:

\noindent Obols:\par\nopagebreak
\noindent\begin{minipage}[b][1.32cm]{\textwidth}\vfill\includePostScript{ObolesGrecques}{0.95}{0.95}\end{minipage}\par

\noindent Drachmas (1 drachma = 6 obols):\par\nopagebreak
\noindent\begin{minipage}[b][1.46cm]{\textwidth}\vfill\includePostScript{DrachmesGrecques}{0.95}{0.95}\end{minipage}\par

\noindent Talents (1 talent = 6,000 drachmas):\par\nopagebreak
\noindent\begin{minipage}[b][1.46cm]{\textwidth}\vfill\includePostScript{TalentsGrecques}{0.95}{0.95}\end{minipage}\par

\noindent For example the number 239,686,966 is represented by:\par\nopagebreak
\noindent\begin{minipage}[b][1.89cm]{\textwidth}\vfill\includePostScript{239686966Grecque}{0.95}{0.95}\end{minipage}\par

\tiret Roman numeration \cite{guitel} is additive and subtractive (any sign placed to the left
of a sign of greater value is subtracted from it). Its digits
are well known:\par\nopagebreak
\noindent\begin{minipage}[b][1.70cm]{\textwidth}\vfill\includePostScript{ChiffresRomains}{0.95}{0.95}\end{minipage}\par

For large numbers, the horizontal bar multiplies by 1,000 the value of the number it
overlines, and the rectangle without a base multiplies by 100,000 the value of the number it encloses,
as in the example 135,792,468 below:\par\nopagebreak
\noindent\begin{minipage}[b][1.90cm]{\textwidth}\vfill\includePostScript{135792468Romain}{0.95}{0.95}\end{minipage}\par

Write a \python\ program that converts a number into Roman
nu\-mer\-a\-tion.

\tiret Write a program that draws the Sierpinski curve of any
order of figure \ref{fig:CourbSierp} (page \pageref{fig:CourbSierp}).

\tiret A magic square of order $n$ is a square array of $n^{2}$ numbers
such that the sum of the numbers of each row, of each column and of each
di\-ag\-o\-nal is the same. Examples are given in figure \ref{fig:Cmagic}
(page \pageref{fig:Cmagic}). Let $C_{i, j}$ be the number in the
$i^{\rm th}$ row and the $j^{\rm th}$ column of the square, counted from
1 and numbered from bottom to top for the rows and from left to
right for the columns. Write a \python\ program that builds
magic squares of odd order $n$ by first placing 1 in
$C_{i, j} = C_{\frac{n + 1}{2}, n}$ and then assigning the numbers 2, 3, \ldots, $n^{2}$
to the elements $C_{i, j}$ obtained by increasing $i$ and $j$ by 1 modulo $n$ each time
for $n - 1$ steps, and by decreasing $j$ by 1 while leaving $i$ unchanged at every
$n^{\rm th}$ step.

\noindent $\Box$}\end{exercices}
\begin{figure}[htbp]
 \includevspacePostScript[9.2cm]{CourbeSierpinski0a5}{0.95}{0.95}
 \caption{Sierpinski curves of orders 0 to 5}\label{fig:CourbSierp}
\end{figure}
\begin{figure}[htbp]
 \includevspacePostScript[6.0cm]{CarreMagique}{0.95}{0.95}
 \caption{Magic squares of orders 3, 5 and 7}\label{fig:Cmagic}
\end{figure}

\pagebreak

\begin{corrige}{
We give one program (without the body of the technical procedure {\idt writeCentred}).
The others are in the folder\newline {\idt book/\allowbreak 22\_variables}:
\input python/MagicSquare.tex
\noindent $\Box$}\end{corrige}
\vfill

%%%%%%%%%%%%%%%%%%%%%%%%%%%%%%%%%%%%%%%%%%%%%%%%%%%%%%%%%%%%%%%%%%%%%%%%%%%
\noindent \framebox[\textwidth]{
\begin{minipage}{\myparboxwidth}
\setlength{\parindent}{2em}
\small \sf
\vspace{\abovedisplayskip}
\section*{Addenda}
A \underline{variable}\index{Variable} is used to give a {\em name} to a memory that keeps
a {\em value} represented
according to a code determined by the {\em type} of the variable. For a binary encoding, one
can imagine that the memory is a row of lamps that are off (to represent
the bit 0) or on (for the bit 1).
In many languages, such as \pascal, the initial
value of a va\-ri\-a\-ble is {\em undetermined}, since it depends on the state
of the memory at the moment the execution of the program begins, which is a source
of errors. In \python, a variable receives an initial value as soon as it is created.
The \underline{assignment}\index{Assignment} instruction makes it possible to
keep a new value by physically mod\-i\-fy\-ing
the state of this memory. The new value is given by an expression (which can
use the old value). The computing notion of variable is therefore very different
from the mathematical notion of the same name. In mathematics a variable denotes an
unknown value fixed once and for all, like the constants of a
\python\ program whose value is never changed. For a mathematician, a computer variable
would be the sequence of values it takes when it is created and then after each
assignment.

\vspace{\baselineskip}


Variables {\em local} to a procedure (or function) can be declared
by writing {\idt x =} $\ldots$ in the body of the procedure, like
the variable {\idt n} of the function {\idt readOrder} of the program {\idt MagicSquare}.
In this case the variable can be used only in the rest of the body of the procedure.
At run time, it is created at each execution of the body
of the procedure. For a recursive procedure,
it is at each call a new variable without any
link to the instances corresponding to earlier recursive calls.

\vspace{\abovedisplayskip}
\end{minipage}
}

%%%%%%%%%%%%%%%%%%%%%%%%%%%%%%%%%%%%%%%%%%%%%%%%%%%%%%%%%%%%%%%%%%%%%%%%%%%%
%%%%%%%%%%%%%%%%%%%%%%%%%%%%%%%%%%%%%%%%%%%%%%%%%%%%%%%%%%%%%%%%%%%%%%%%%%%
\noindent \framebox[\textwidth]{
\begin{minipage}{\myparboxwidth}
\setlength{\parindent}{2em}
\small \sf
\vspace{\abovedisplayskip}
In the body of the procedure one can also use
{\em glo\-bal} variables whose declaration appears in the program,
outside the procedures, such as {\idt width} and {\idt height} in the program
{\idt ComputeRobotWindowSize} or {\idt i} and {\idt j} in the program {\idt MagicSquare}.
For a procedure to be able to modify them by assignment, they must be named after the
keyword {\mcl global} at the beginning of its body. These global variables exist
throughout the execution of the program. Consequently, in the case of
recursive calls of a procedure, it is always the
same global variable.

Finally, local and global variables can have the same name (like the variable
{\idt n} of the function {\idt readOrder} and the variable {\idt n} of the program {\idt MagicSquare}). In
this case it is the local variable that is visible in the procedure, the global
variable continuing to exist but being inaccessible. In \python, a variable is local
to a procedure as soon as it receives a value by assignment in the body of this procedure
(unless it is declared {\mcl global}).

\vspace{\abovedisplayskip}
\end{minipage}
}
%%%%%%%%%%%%%%%%%%%%%%%%%%%%%%%%%%%%%%%%%%%%%%%%%%%%%%%%%%%%%%%%%%%%%%%%%%%%

